\section{Implementation}
\label{sec:implementation}

Table~\ref{tab:stack} summarizes the technology stack.
The open-source repository is available at \url{https://github.com/Fatihmaull/halal-chain}.

\begin{table}[t]
  \caption{HalalChain implementation stack.}
  \label{tab:stack}
  \centering
  \begin{tabular}{@{}ll@{}}
    \toprule
    \textbf{Layer} & \textbf{Technology} \\
    \midrule
    Blockchain & Solidity 0.8.20, OpenZeppelin, Base Sepolia \\
    Development & Hardhat 3, Mocha/Chai (7 tests) \\
    Frontend & Next.js 16, Wagmi v3, Viem, Tailwind \\
    Storage & IPFS via Pinata (server-side upload API) \\
    Deployment & Hardhat scripts, Vercel (frontend) \\
    \bottomrule
  \end{tabular}
\end{table}

\subsection{Smart Contracts}
The contract implements RBAC with \texttt{AccessControl}, replacing a single-owner auditor mapping.
Unit tests cover registration, verification, rejection, revision, revocation, unauthorized access, and invalid batch identifiers.

\subsection{Web Application}
Three dashboards serve distinct roles:
\begin{itemize}
  \item \textbf{Producer} --- multi-step registration, IPFS upload, QR code generation, batch history
  \item \textbf{Auditor} --- pending queue, document review links, verify/reject/revoke with confirmation
  \item \textbf{Consumer} --- \texttt{/verify/\{batchId\}} read-only page; bilingual EN/ID; no wallet required
\end{itemize}

\subsection{IPFS Upload}
A Next.js API route (\texttt{POST /api/ipfs/upload}) pins files via Pinata JWT server-side.
A development mock mode generates deterministic local CIDs when Pinata is not configured; benchmarks use real Pinata pins only (Section~\ref{sec:eval-b}).
